\section{Overview}
\label{sec:overview}

This document supports the major revision of
\emph{``Online activity prediction via generalized Indian buffet
process models,''} submitted to the Annals of Applied Statistics
(AOAS), following reports from Editor Po-Ling Loh, an Associate
Editor (AE), Reviewer~1 (R1), and Reviewer~2 (R2).

It serves two purposes:
\begin{enumerate}[leftmargin=*]
  \item \textbf{Technical reference.}
        A self-contained record of the data pipeline, experiment
        construction, and evaluation protocol for the new real-data
        analyses added in the revision.
  \item \textbf{Rebuttal scaffold.}
        The problem--solution structure (Section~\ref{subsec:problems})
        provides the narrative arc of the response letter, while the
        per-reviewer item tracker (Section~\ref{subsec:themes})
        provides the skeleton for the point-by-point response.
\end{enumerate}


% =========================================================
\subsection{Reviewer Reports}
\label{subsec:reviewer_reports}
% =========================================================

We summarize the reports below; the full verbatim text is
reproduced in Appendix~\ref{app:reviews}.

\paragraph{Editor and AE.}
The AE recommends: \emph{``focus more on the data analysis and the
downstream impact.''}
The editor requests a dedicated Section~2 for the motivating
application and dataset, a $\leq$20-page manuscript, minimal color
in figures, and a single point-by-point response file.

\paragraph{Reviewer~1.}
Positive overall (\emph{``I thoroughly enjoyed reading this paper''}).
Eleven specific points, mostly notation and clarity.
The most substantive (R1-1) requests a rewrite of Section~3.4
to clarify $D_M$, ``new'' vs.\ ``unobserved'' users, and the
relationship between $F'$ and $Z^{\mathrm{FT}}$.

\paragraph{Reviewer~2.}
The critical reviewer. Four major concerns:
(a)~novelty vs.\ Camerlenghi et al.\ (2022);
(b)~paper reads as methodology, data analysis is thin,
\emph{``little evidence the method has practical impact''};
(c)~primary goal unclear (causal estimation vs.\ experimental design);
(d)~theoretical results lack practical interpretation.
Plus 12 minor points and several typos.


% =========================================================
\subsection{Problems and Solutions}
\label{subsec:problems}
% =========================================================

We distill the reports into five problems.
Each problem statement can be adapted directly into the
rebuttal letter; each solution summarizes the revision action
and points to the section of this document (and the revised
manuscript) where it is carried out.

\paragraph{Problem 1: The data analysis is thin.}
\begin{itemize}[leftmargin=*]
  \item[\emph{Source:}] AE; R2 major~(b); R2 minor~(j).
  \item[\emph{Problem:}]
        The original submission relied on a proprietary dataset
        (1,774 A/B test arms, not shareable) and the ASOS dataset
        (76 treatment arms, first-trigger counts only).
        Neither was described in sufficient detail, and the analyses
        were limited to accuracy comparisons without diagnostics.
  \item[\emph{Solution:}]
        We introduce two new public datasets---REES46 (380M events,
        7~months, per-user timestamps) and UCI Online Retail
        (541K transactions, 374~days).
        These provide the fine-grained structure needed to evaluate
        all three model variants and yield 20+ experiment windows
        for systematic evaluation.
  \item[\emph{Where:}]
        Data sources: Section~\ref{sec:data_sources}.
        Preprocessing: Section~\ref{sec:preprocessing}.
        Manuscript: revised Section~2 and Section~6.
\end{itemize}

\paragraph{Problem 2: The practical impact is unclear.}
\begin{itemize}[leftmargin=*]
  \item[\emph{Source:}] AE; R2 major~(b), (c).
  \item[\emph{Problem:}]
        The reviewers found \emph{``little evidence the method has
        practical impact and changes decisions.''}
        The paper did not demonstrate how predictions translate into
        actionable experiment-duration decisions.
  \item[\emph{Solution:}]
        We commit to the framing: the method supports
        \emph{experiment-duration planning}.
        We demonstrate impact through:
        (i)~calibration diagnostics (95\% CI coverage across all
        experiment windows);
        (ii)~case studies showing early-stopping decisions on 3
        representative experiments;
        (iii)~$D_M$ interval computations answering ``how many more
        days to reach $M$ users?''
  \item[\emph{Where:}]
        Evaluation protocol: Section~\ref{sec:fitting}.
        Manuscript: revised Section~6 and new Section~2 framing.
\end{itemize}

\paragraph{Problem 3: No evidence for the power-law claim.}
\begin{itemize}[leftmargin=*]
  \item[\emph{Source:}] R2 minor~(d).
  \item[\emph{Problem:}]
        The paper asserted that user activity exhibits power-law
        growth without providing a reference or visualization.
  \item[\emph{Solution:}]
        Log-log histograms of per-user trigger counts and cumulative
        new-user curves for both datasets, paired with a citation
        to~\citet{clauset2009power}.
  \item[\emph{Where:}]
        Section~\ref{sec:preprocessing}.
        Manuscript: new Figure~1 in Section~2.
\end{itemize}

\paragraph{Problem 4: The paper's primary goal is unclear.}
\begin{itemize}[leftmargin=*]
  \item[\emph{Source:}] R2 major~(c); AE.
  \item[\emph{Problem:}]
        The introduction oscillated between improving causal effect
        estimation and improving experimental design.
  \item[\emph{Solution:}]
        Single framing: the goal is \emph{prediction of user
        engagement to inform experiment duration decisions}.
        Estimands: $\Utrue$, $\Ttrue$, $D_M$.
        This is a design/planning contribution, not causal estimation.
  \item[\emph{Where:}]
        Manuscript: rewritten introduction and new Section~2.
\end{itemize}

\paragraph{Problem 5: NB-SSP never evaluated on public data.}
\begin{itemize}[leftmargin=*]
  \item[\emph{Source:}] R2 major~(b).
  \item[\emph{Problem:}]
        The \nbmodel{} model's total-trigger predictions ($\Ttrue$)
        were only evaluated on proprietary data.
        The ASOS dataset provides only first-trigger counts.
  \item[\emph{Solution:}]
        Both REES46 and UCI provide per-user daily trigger counts,
        enabling the first public-data evaluation of \nbmodel{}'s
        total-trigger predictions.
  \item[\emph{Where:}]
        Section~\ref{sec:fitting}.
        Manuscript: new results in Section~6.
\end{itemize}


% =========================================================
\subsection{Per-Reviewer Item Tracker}
\label{subsec:themes}
% =========================================================

For the point-by-point rebuttal, we track every individual
reviewer item organized by theme.
Table~\ref{tab:themes} gives the summary; the full item-level
detail is in the companion files
\texttt{t1\_framing.tex} through \texttt{t6\_typos.tex}.

\begin{table}[H]
\centering
\caption{Revision themes and status.
Themes T1--T2 are addressed by the data pipeline documented
in this appendix (Sections~\ref{sec:data_sources}--\ref{sec:fitting}).
Themes T3--T6 are addressed in the revised manuscript text.}
\label{tab:themes}
\small
\begin{tabular}{@{}clp{3.2cm}cc@{}}
\toprule
\textbf{ID} & \textbf{Theme} & \textbf{Primary source} &
\textbf{Open} & \textbf{Done} \\
\midrule
T1 & Framing \& Impact & AE, R2-b,c, Editor & 4 & 0 \\
T2 & Real Data Analysis & AE, R2-b,j, Editor & 4 & 0 \\
T3 & Theory & R2-a,d,f--h,l & 1 & 9 \\
T4 & Sec.~3.4 \& Notation & R1-1 through R1-11 & 11 & 0 \\
T5 & Simulations & R2-k, Editor & 3 & 0 \\
T6 & Typos & R1, R2, proofread & 16 & 0 \\
\midrule
   & \textbf{Total} & & \textbf{39} & \textbf{9} \\
\bottomrule
\end{tabular}
\end{table}

Table~\ref{tab:problem_theme_map} cross-references the five
problems with the theme items they resolve, so that each
problem--solution pair in the rebuttal can cite the specific
items it addresses.

\begin{table}[H]
\centering
\caption{Cross-reference: problems $\leftrightarrow$ theme items.}
\label{tab:problem_theme_map}
\small
\begin{tabular}{@{}clll@{}}
\toprule
\textbf{Problem} & \textbf{Short title} &
\textbf{Theme items} & \textbf{Appendix section} \\
\midrule
1 & Thin data analysis &
  T1.d, T2.a, T2.b, T2.c &
  Secs.~\ref{sec:data_sources}--\ref{sec:preprocessing} \\
2 & Practical impact &
  T1.a, T1.b, T1.c, T2.e &
  Sec.~\ref{sec:fitting} \\
3 & Power-law evidence &
  T2.b, T2.c &
  Sec.~\ref{sec:preprocessing} \\
4 & Unclear primary goal &
  T1.a &
  Manuscript intro + Sec.~2 \\
5 & NB-SSP on public data &
  T2.a, T2.e &
  Sec.~\ref{sec:fitting} \\
\bottomrule
\end{tabular}
\end{table}


% =========================================================
\subsection{Evaluation Strategy}
\label{subsec:eval_strategy}
% =========================================================

The solutions above require four concrete evaluation exercises.
Table~\ref{tab:eval_exercises} lists them; each will be cited
in the rebuttal when responding to the corresponding reviewer
item.

\begin{table}[H]
\centering
\caption{Evaluation exercises.
Each exercise produces results cited in the rebuttal and
incorporated into the revised Section~6.}
\label{tab:eval_exercises}
\small
\begin{tabular}{@{}clp{4cm}l@{}}
\toprule
\textbf{ID} & \textbf{Use case} &
\textbf{What we compute} & \textbf{Problem} \\
\midrule
E1 & Accuracy &
  $\hat{U}$ + 95\% CI for all windows;
  accuracy ratio $v$ &
  P1, P2 \\
E2 & Stopping decisions &
  $D_M$ intervals for 3 experiments &
  P2 \\
E3 & Trigger forecasting &
  $\hat{T}$ via \nbmodel{} &
  P5 \\
E4 & Competitor comparison &
  E1/E3 for BB, HBG, Jackknife, GT, LP &
  P1, P2 \\
\bottomrule
\end{tabular}
\end{table}

\paragraph{Two layers of evidence.}
Results are presented at two levels:
\begin{enumerate}[leftmargin=*]
  \item \textbf{Aggregate:} accuracy boxplots, CI coverage tables,
        and trigger-prediction accuracy across all experiment windows.
  \item \textbf{Case studies:} prediction trajectories, $D_M$
        intervals, and a decision narrative for 3 representative
        experiments (sparse, medium, high traffic).
\end{enumerate}


% =========================================================
\subsection{Document Roadmap}
\label{subsec:roadmap}
% =========================================================

The remainder of this appendix follows the data pipeline
from raw files to model evaluation:
\begin{description}[leftmargin=!,labelwidth=2.8cm]
  \item[Section~\ref{sec:data_sources}]
    Data sources: REES46 and UCI Online Retail
    (addresses P1, P5).
  \item[Section~\ref{sec:preprocessing}]
    Preprocessing: chunked loading, aggregation,
    user-engagement characterization
    (addresses P1, P3).
  \item[Section~\ref{sec:experiment_construction}]
    Experiment construction: windowing strategies,
    ground-truth definitions.
  \item[Section~\ref{sec:fitting}]
    Model fitting and evaluation protocol: exercises
    E1--E4 (addresses P2, P5).
  \item[Appendix~\ref{app:reviews}]
    Full verbatim reviewer reports.
\end{description}
